\documentclass{article}%
\usepackage[T1]{fontenc}%
\usepackage[utf8]{inputenc}%
\usepackage{lmodern}%
\usepackage{textcomp}%
\usepackage{lastpage}%
\usepackage{expex}%
\usepackage{geometry}%
\geometry{margin=2cm}%
\usepackage[spanish]{babel}%
%
\lingset{everygla=\it}%
\title{Análisis Lingüístico: Gramatica Normativa Kaqchikel}%
\author{Generado por el TFM de Elías Carrasco}%
%
\begin{document}%
\normalsize%
\maketitle%
\section{Page 203}%
\label{sec:Page203}%

%
\subsection{Oración compleja}%
\label{subsec:Oracincompleja}%

%
La oración compleja se forma con más de una cláusula, una de ellas es la principal y %
otra la subordinada. La idea que expresan las cláusulas subordinadas siempre dependen %
de la cláusula principal, su posición es siempre después de ésta. Una cláusula subordinada %
se introduce por medio de subordinadores, aunque hay excepciones. Son oraciones %
complejas si dentro de la cláusula principal se encuentran las siguientes cláusulas %
subordinadas: %
\subsection{Cláusula relativa}%
\label{subsec:Clusularelativa}%

%
\subsection{Cláusula adjunta o adverbial}%
\label{subsec:Clusulaadjuntaoadverbial}%

%
\subsection{Cláusula de complemento}%
\label{subsec:Clusuladecomplemento}%

%
\subsection{Cláusula relativa}%
\label{subsec:Clusularelativa}%

%
La cláusula relativa modifica a sustantivos, describe alguna característica, por eso %
funciona como adjetivo. La cláusula relativa se introduce por los subordinadores: ri, %
re, la, achoq rik'in y achoq chi re o simplemente sin subordinador. Estas partículas se %
conocen con los nombres de pronombre relativo o relativizador. Si modifica al sujeto u %
objeto, se introduce con las partículas: ri, re, la. Si modifica al instrumento o compañía, %
la cláusula relativa se introducen con la frase achoq rik'in; y si la cláusula no se %
introduce por subordinador, la cláusula relativa se escribe entre comas. Ejemplos: %
\subsection{Modificando al sujeto:}%
\label{subsec:Modificandoalsujeto}%

%
\ex
  	abla Ri ixöq, ri xkos, xb'eruloq'o' xkoya'.
  	ablz Ri ixöq ri x- 0- kos x- 0- b'e-ru-loq'-o' xkoya'.
  	able \textbf{DET} \textbf{mujer} \textbf{CR} \textbf{COM-} \textbf{B3s-} \textbf{cansar} \textbf{COM-} \textbf{B3s-} \textbf{MOV-A3s-comprar-SC} \textbf{tomate}
  	ablf [Ri ixöq]\textsubscript{CR}
  \xe[\textit{La mujer que se cansó fue a comprar tomate}]
%
\ex
  	abla Re ak'wal, re xtzaq, xoq'.
  	ablz Re ak'wal re x- 0- tzaq xoq'.
  	able \textbf{DET} \textbf{niño} \textbf{CR} \textbf{COM-} \textbf{B3s-} \textbf{caer} \textbf{llorar}
  	ablf [Re ak'wal]\textsubscript{CR}
  \xe[\textit{Este niño que se cayó, lloró.}]
%
\ex
  	abla La k'ajöl, la xloq'o ri wexaj, xb'e.
  	ablz La k'ajöl la x- loq'o ri wexaj xb'e
  	able \textbf{DET} \textbf{muchacho} \textbf{CR} \textbf{COM-} \textbf{comprar} \textbf{DET} \textbf{pantalón} \textbf{se fue}
  	ablf [La k'ajöl]\textsubscript{CR}
  \xe[\textit{Ese muchacho que compró el pantalón, se fue.}]
%
\end{document}